%% Chapter 9 -----------------------------------------------------------------
\chapter{Type II Band-Splitting Analysis}
\label{ch:band-splitting}
\index{band splitting}

The emission lanes of many Type~II bursts are split into two parallel bands, interpreted
as emission from the plasma ahead of (upper band) and behind (lower band) the shock front.
The frequency ratio of the two bands gives the shock's density compression ratio\index{compression ratio}, from
which the Alfvén Mach number\index{Alfvén Mach number}, the Alfvén speed\index{Alfvén speed} and the coronal magnetic field\index{magnetic field} can be
estimated \parencite{vrsnak2002}. The Type~II Band-splitting window performs this analysis
directly on the current noise-reduced dynamic spectrum.

\section{Opening the window}

Open the window with \menu{Analysis > Type II Band-splitting > Open Type II
Band-splitting}. The calculation needs the shock speed and the starting frequency from the
Burst Analyzer, so run the Analyzer (Section~\ref{sec:analyzer}) first. The
\gui{Analyzer Reference}\index{band splitting!Analyzer reference} panel reports \gui{Analyzer reference data loaded.} when these
values are available, and \gui{Analyzer input missing.} otherwise.

\section{Window layout}

The tools and panels of the window are numbered in Figure~\ref{fig:band-splitting} and
described in Table~\ref{tab:band-splitting-tools}.

\screenshot{analysis/typeII_band_splitting_window}{The Type~II Band-splitting window with its tools and panels numbered.}{fig:band-splitting}{Band-splitting window with numbered callouts.}

\begin{table}[!htbp]
  \centering
  \caption{Tools and panels of the Type~II Band-splitting window.}
  \label{tab:band-splitting-tools}
  \small
  \begin{tabularx}{\linewidth}{@{}cP{0.25\linewidth}L@{}}
    \toprule
    \thead{No.} & \thead{Tool or panel} & \thead{Function} \\
    \midrule
    \callout{1} & Add Points & Toggle point picking; click along the active band on the spectrum to add points. \\
    \callout{2} & Undo Last Point & Remove the last point added to the active band. \\
    \callout{3} & Clear Active Band & Remove all points of the active band. \\
    \callout{4} & Fit Active Band & Fit the points of the active band with a power law. \\
    \callout{5} & Fit Both Bands & Fit the upper and the lower band. \\
    \callout{6} & Save Plot & Save the fitted lanes over the spectrum, or the magnetic field versus height plot, as a figure. \\
    \callout{7} & BvR & Toggle the plot of magnetic field versus shock height (Section~\ref{sec:bvr}). \\
    \callout{8} & Settings & Graph settings: fonts, lane colors, line widths and marker sizes (Section~\ref{sec:bs-graph-settings}). \\
    \callout{9} & Controls panel & \gui{Active Band} (\gui{Upper Band} or \gui{Lower Band}), \gui{Shock Speed} (\gui{Initial Shock Speed}\index{shock parameters!initial speed} or \gui{Average Shock Speed}\index{shock parameters!average speed}) and \gui{Calculate}. \\
    \callout{10} & Results panels & \gui{Analyzer Reference}, \gui{Fitted Bands} and \gui{Averaged Band-Splitting Quantities} (Section~\ref{sec:bs-results}). \\
    \bottomrule
  \end{tabularx}
\end{table}

\section{Picking and fitting the bands}
\label{sec:bs-fitting}

\begin{steps}
  \item Select the \gui{Active Band} in the Controls panel \callout{9}.
  \item Switch on \gui{Add Points} \callout{1} and click points along that band on the
    spectrum. Points of the upper and lower bands are drawn in different colors.
  \item Select the other band and add its points.
  \item Fit each band with \gui{Fit Active Band} \callout{4}, or both at once with
    \gui{Fit Both Bands}\index{band splitting!fitting} \callout{5}.
\end{steps}
Each band is fitted with a power law \(f(t) = a\,t^{-b}\), where \(t\) is measured from the
start of the file. At least two points with positive time and frequency are needed per
band. The \gui{Fitted Bands} panel shows each fitted equation, \(R^2\), RMSE and the number
of selected points.

\section{Calculation and results}
\label{sec:bs-results}

Choose whether the \gui{Initial Shock Speed} or the \gui{Average Shock Speed} from the
Analyzer is used, and click \gui{Calculate}. The results panels contain:
\begin{description}
  \item[Analyzer Reference] the reference values imported from the Burst Analyzer: the
    fold number\index{fold number}, the starting frequency \(f_\mathrm{start}\), the average drift rate
    \(\langle\mathrm{d}f/\mathrm{d}t\rangle\), the initial and average shock speeds
    \(V_{s,0}\) and \(\langle V_s\rangle\), and the initial and average shock heights
    \(R_{p,0}\) and \(\langle R_p\rangle\).
  \item[Fitted Bands] the fitted equation, \(R^2\), RMSE and number of points of each band.
  \item[Averaged Band-Splitting Quantities] the averaging interval, the start time
    \(t_\mathrm{start}\) and band frequencies, the mean upper and lower frequencies, the
    bandwidth \(\langle\Delta f\rangle = \langle f_u - f_l\rangle\), the mean upper-band
    drift \(\langle\mathrm{d}f_u/\mathrm{d}t\rangle\), the compression ratio
    \(X = (\langle f_u\rangle/\langle f_l\rangle)^2\), the Alfvén Mach number \(M_A\), the
    Alfvén speed \(V_A\) and the magnetic field \(B\).
\end{description}
The averaging interval is the time span of the upper-band points; the two fitted curves are
sampled over it. If the lower-band fit has to be extrapolated over part of the interval, a
warning is shown. The calculation is refused if the compression ratio does not stay
between 0 and 4 over the interval.

\begin{background}[Background: band-splitting relations]
  For a perpendicular shock in a low-\(\beta\) plasma the compression ratio \(X\) and the
  Alfvén Mach number are related by \parencite{vrsnak2002}
  \begin{equation*}
    X = \left(\frac{f_u}{f_l}\right)^2, \qquad
    M_A = \sqrt{\frac{X\,(X+5)}{2\,(4-X)}}.
  \end{equation*}
  The Alfvén speed follows from the shock speed \(V_s\) taken from the Analyzer,
  \(V_A = V_s/M_A\), and the magnetic field from
  \(B = V_A\sqrt{\mu_0\,m_p\,n_e}\), with the electron density obtained from the plasma
  frequency, \(n_e = (f/8.98\times10^{-3}\,\mathrm{MHz})^2\)~cm\(^{-3}\). In the averaged
  result the density is taken at the Analyzer's starting frequency.
\end{background}

\section{Magnetic field versus height}
\label{sec:bvr}
\index{magnetic field!versus height}

After a successful calculation, the \gui{BvR}\index{band splitting!B vs R plot} tool \callout{7} plots the magnetic field
against shock height (Figure~\ref{fig:bvr}). Along the averaging interval the field is
computed at each sample time from the local compression ratio, with the density taken from
the upper-band frequency. The heights are calculated with the Newkirk\index{density model!Newkirk} model and the
Analyzer's fold number, and the profile is fitted with a power law \(B = a\,R^{-b}\). Click
\gui{BvR} again to return to the spectrum.

\screenshot{analysis/BvR}{Magnetic field versus shock height in the band-splitting window.}{fig:bvr}{Band-splitting window showing the B vs R plot.}

\begin{note}
  The shock speed is taken from the Burst Analyzer and therefore follows the density model
  selected there (Section~\ref{sec:density-models}). The heights of the magnetic field
  versus height profile always use the Newkirk model, with the fold number set in the
  Analyzer.
\end{note}

\section{Graph settings}
\label{sec:bs-graph-settings}

The \gui{Settings} tool \callout{8} opens the \gui{Type II Graph Settings}\index{band splitting!graph settings} dialog
(Figure~\ref{fig:bs-settings}):
\begin{description}
  \item[Text Formatting] font family, and size, bold and italic for the title, the axis
    labels and the tick labels.
  \item[Upper Band, Lower Band] color and width of the fit line, and color and size of
    the picked points.
  \item[B vs R] color and width of the fit line, and color and size of the markers.
\end{description}
\gui{Apply} applies the settings, \gui{Reset to Defaults} restores the defaults and
\gui{Close} closes the dialog.

\screenshot[0.45\linewidth]{dialogs/type_ii_graph_settings}{The \gui{Type II Graph Settings} dialog.}{fig:bs-settings}{The Type II Graph Settings dialog showing the Text Formatting, Upper Band, Lower Band and B vs R groups.}

\section{Saving plots}

\gui{Save Plot} \callout{6} exports either the fitted lanes over the dynamic spectrum or the
magnetic field versus height fit, whichever is on show, as an OriginPro-style graph in PNG,
PDF, EPS, SVG, TIFF or JPG format. The lane colors and fonts from the graph settings are
kept.

\begin{caution}
  Derived magnetic-field values may not be accurate for all events or under all
  assumptions. They depend on the shock geometry, the density model and the
  interpretation of the split bands. Confirm results against known or independently
  validated events before using them for scientific conclusions.
\end{caution}
